% =====================================================================
% Section 3: the case forall x phi (key: univ). Source: research/tracks/universal/notes-final.md
% (and its referee report), with E1/E3(a) from the experiments track. Proofs in app-universal.tex.
% =====================================================================

\section{The case \texorpdfstring{$\forall x\varphi$}{forall x phi}: what instances support}\label{sec:univ}

H\"anni's question starts from one case, statements $\varphi(t)$ and the law $\forall x\varphi$: ``forall x phi(x) is a simple law which implies all the other statements. so once we see all these other statements, we should up the probability on forall x phi(x). i guess we should still be entertaining other hypotheses as well'' (\Cref{sec:intro}). In the model of \Cref{sec:model} the answer is this. The data favour the generator that emits exactly the instances, and never favour $\forall x\varphi$ over its instance schema $\sigma_\varphi$. ``All future data are instances'' is confirmed; ``the axioms prove $\forall x\varphi$'' is not, a Bayesian form of the $\omega$-gap of AS~\S4.5. $\forall x\varphi$ is licensed by an instance at a parameter, by a quantified datum, or by a likelihood that models which theorems are reported. All results are from the universal track; proofs and further tables are in \Cref{app:univ}. Logarithms are natural; rates are in nats per datum.

% ---------------------------------------------------------------------
\subsection{Setup}\label{sec:univ:setup}

\begin{definition}[Instances, schema, term laws]\label{def:univ:setup}\src{universal \S1.1--1.3}
$\varphi(x)$ is an $L_A$-formula with exactly one free variable $x$, which occurs in it; the examples ($0+x=x$, $x+0=x$, $x\cdot0=0$, $\neg Sx=0$, $x<Sx$) are quantifier-free unless stated. $\Ic(\varphi):=\{\varphi(t):t\text{ closed}\}$; $\Io(\varphi)$ also allows parameters in $t$. The instance schema is $\sigma_\varphi:=\varphi[z/x]$, $z$ a $0$-ary term metavariable, closed guard unless stated. Term laws: $\Qnum(S^j0)=(1-q)q^j$, $q=\tfrac12$ by default; $\QGW$, each node picking its root $0,S,+,\cdot$ with probabilities $.4,.3,.2,.1$; $\Qopen:=\rho[p]+(1-\rho)\Qnum$ for a parameter $p$. The first two are root-generated: $\Qg(f(t_1,\dots,t_r))=p_f\prod_i\Qg(t_i)$. One law supplies every metavariable body and every $\forall$E term.
\end{definition}

Calculi and likelihoods are as in \Cref{def:model:calculi,def:model:variants}. In $\Cmin$ each chain node cites (probability $1-c$) or applies $\forall$E with $t\sim\Qg$ (probability $c$); $\Copen(c,g)$ adds Gen on the parameter ($g$; cite $p_{\mathrm{cite}}=1-c-g$); $\mu_T$ is the unnormalised derivation mass and $\Lone=\mu_T/Z_T$. H\"anni's scores are $\Sprove$, $\Snc$ and $\Sg$ with $g\equiv1$ on provable data (the track's $S_{\mathrm{nc},\beta}$ is $\beta^n\Sg$ with $g_\infty=1/\beta$; \Cref{def:model:scores}). Theorems about countable classes assume a proper prior (\Cref{def:model:prior}); the computations use finite or summable families with the track's prior $2^{-(\sum_A\lvert A\rvert+\#\text{axioms})}$, improper over all theories, and memorisers with a product-Bernoulli prior over a named universe ($r_s=2^{-(\lvert s\rvert+1)}$, $\sum r_s<\infty$). Defaults: $c=0.3$; $(c,g,\rho,q)=(0.3,0.2,0.1,\tfrac12)$ in $\Copen$; Dirichlet $\alpha=1$ in checks c2--c7, $\alpha=\tfrac12$ in c8--c10.

\begin{table}[t]
\centering\footnotesize
\begin{tabularx}{\textwidth}{@{}l>{\raggedright\arraybackslash}X c@{}}
\toprule
name & axioms & $\vdash\forall x\varphi$: pure / over $\Q$\\
\midrule
$\Hall$; $\Hsch$ & $\{\forall x\varphi\}$; $\{\sigma_\varphi\}$, closed guard & yes/yes; no/no\\
$\Hopen$; $\Hboth$ & $\{\sigma_\varphi\}$, open guard ($z\sim\Qopen$); $\{\forall x\varphi,\sigma_\varphi\}$, weight $w_\forall$ on $\forall x\varphi$ & yes/yes; yes/yes\\
memorisers; over-specific & finite sets of closed instances; e.g.\ $\sigma_\varphi[z:=Sz]$, closed guard & no/no; no/no\\
over-general & a subterm or subformula of $\sigma_\varphi$ replaced by a fresh metavariable; the bare $F_0$ (inconsistent) & varies\\
root split & $\{\sigma_\varphi[z:=f(z_1,\dots,z_r)]:f\text{ a root of }\Qg\}$, closed guards & no/no\\
$R_k$ & $\{\varphi(0),\dots,\varphi(S^{k-1}0),\varphi(S^kz)\}$, $z$ unguarded; weights $(v_0,\dots,v_{k-1},u)$ & no/yes\\
$\Split_k$; Both0 & $\{\varphi(0),\dots,\varphi(S^{k-1}0),\forall y\varphi(S^ky)\}$; $\{\varphi(0),\forall x\varphi\}$ (sentences only) & no/yes; yes/yes\\
\bottomrule
\end{tabularx}
\caption{Hypotheses of \Cref{sec:univ}\src{universal \S1.3}; provability for $\varphi=0+x=x$. A theory of closed instances proves $\forall x\varphi$ neither in pure logic nor over $\Q$ (AS Ex~4.9); $R_k$, $\Split_k$: \Cref{prop:univ:noguard,prop:univ:sentences}. $B\oplus_wA$ adds $A$ with weight $w$ to $B$, scaling $B$'s weights by $1-w$.}
\label{tab:univ:hyp}
\end{table}

% ---------------------------------------------------------------------
\subsection{Instance data never favour \texorpdfstring{$\forall x\varphi$}{forall x phi} over its schema}\label{sec:univ:never}

\begin{proposition}[Factorisation]\label{prop:univ:factor}\status{proved; computed}\src{universal Prop U1}
In $\Cmin$, let $\varphi$ have one free variable and $m$ leading quantifiers. Then $\mu_{\Hall}(s)=c\,\mu_{\Hsch}(s)$ for every sentence $s\neq\forall x\varphi$; $\mu_{\Hall}(\forall x\varphi)=1-c$, $\mu_{\Hsch}(\forall x\varphi)=0$; $Z_{\Hsch}=(1-c)\sum_{k\le m}c^k$, $Z_{\Hall}=(1-c)+cZ_{\Hsch}$. With a background, $\mu_{B\oplus_w\forall x\varphi}=(1-w)\mu_B+wc\,\mu_{\Hsch}$ and $\mu_{B\oplus_w\sigma_\varphi}=(1-w)\mu_B+w\,\mu_{\Hsch}$ on $s\neq\forall x\varphi$. For $k$ variables the factor is $c^k$ on each full instance.
\end{proposition}

\noindent\emph{Proof idea.} ``Cite $\forall x\varphi$, eliminate with $t_1,\dots,t_k$'' maps to ``cite $\sigma_\varphi$ with $z:=t_1$, eliminate with $t_2,\dots,t_k$'': same conclusion, probability larger by $1/c$, a bijection onto all $\sigma_\varphi$-chains (\Cref{app:univ:odds}). Checked to $5\cdot10^{-14}$ on $2\cdot10^5$ sampled derivations.

\begin{theorem}[Odds on instance data]\label{thm:univ:odds}\status{proved; computed}\src{universal Thm U2}
In $\Cmin$, after $n$ closed instances the odds $O_n:=\pi_n(\Hall)/\pi_n(\Hsch)$ satisfy $O_n/O_0$ as in \Cref{tab:univ:odds}.
\end{theorem}

\begin{table}[t]
\centering\small
\begin{tabularx}{\textwidth}{@{}>{\raggedright\arraybackslash}p{0.47\textwidth}X@{}}
\toprule
likelihood & $O_n/O_0$\\
\midrule
$\Lzero$ (strict citation) & $0$ for $n\ge1$\\
$\Lcl$; $\Lsel(S)$ with $\forall x\varphi\notin S$; $\Sprove$, $\Snc$, $\Sg$ ($\forall x\varphi$ consistent) & $1$\\
$\mu_T$ (unnormalised); $\Lmax$ & $c^n$\\
$\Lone$ & $\big(cZ_{\Hsch}/((1-c)+cZ_{\Hsch})\big)^n$; $(c/(1+c))^n$ for quantifier-free $\varphi$\\
background $B$, under $\mu_T$ & $\prod_ir(X_i)$, $r(s)\in[c,1]$; $r(s)=c$ if $\mu_B(s)=0$\\
\bottomrule
\end{tabularx}
\caption{Odds of $\Hall$ against $\Hsch$ after $n$ closed instances, relative to the prior odds (\Cref{thm:univ:odds}); asserted on every sampled path (5 formulas, 2 term laws, 50 seeds, $n\le200$), referee's code to $7\cdot10^{-13}$. The consistency condition in the score row is ours; the notes leave it implicit.}
\label{tab:univ:odds}
\end{table}

\begin{theorem}[Instance data never favour $\forall x\varphi$]\label{thm:univ:B}\status{proved}\src{universal Thm B}
On data not containing $\forall x\varphi$, the posterior odds of the $\forall$-version against the schema version are at most the prior odds: \textup{(B1)} in $\Cmin$, single axioms, closed instances, under $\Lcl$, $\mu_T$, $\Lone$, $\Lsel$, $\Lmax$ and the scores (equality under $\Lcl$, $\Lsel$, scores); \textup{(B2)} in $\Cmin$, $B\oplus_w\forall x\varphi$ against $B\oplus_w\sigma_\varphi$ for any background $B$ and any data, under $\mu_T$ and $\Lone$, with a fixed $w$ or a common weight prior (equality under $\Lcl$); \textup{(B3)} in $\CUiii$ calculi, quantifier-free $\varphi$, under $\mu_T$, $\Lmax$ and the scores.
\end{theorem}

\noindent\emph{Proof idea.} (B2) holds at every $w$, and $Z_{\Hall}-Z_{\Hsch}=(1-c)(1-Z_{\Hsch})\ge0$ orders the normalisers. (B3) rests on a simulation inequality: in $\CUiii$ every citation of $\forall x\varphi$ sits under a $\forall$E, and replacing the pair by a citation of $\sigma_\varphi$ divides the probability by $c$ (\Cref{prop:univ:sim}). The mechanism is the size principle (\Cref{lem:model:size}): $\forall x\varphi$ spends mass on outputs never observed. Not covered: $\Lsel$ with a background and learned weights (\Cref{prop:univ:B2}), $\wedge$-detours (\Cref{rem:univ:detour}), normalised likelihoods in $\CUiii$.

\begin{proposition}[Robustness of the direction]\label{prop:univ:robust}\status{proved; computed}\src{universal Props U2n, U2t, U2c}
In $\Cmin$, for every $s\neq\forall x\varphi$: \textup{(U2n)} with a full-support noise law $N$ and $\eta\in(0,1)$, the ratio of the noisy laws of $\Hall$ and $\Hsch$ at $s$ is in $[c,1]$ under $\mu_T$, $\le1$ under $\Lone$, and $1$ under $\Lsel$ with noise after selection; a false datum no longer refutes. \textup{(U2t)} If chains are weighted by $f(\ell)w_A\prod_j\Qg(t_j)$, $f$ non-increasing in the number $\ell$ of eliminations (e.g.\ Levin-style $f(\ell)\propto(\ell+1)^{-2}$), then $P^f_{\Hall}(s)\le\sup_\ell\frac{f(\ell+1)}{f(\ell)}P^f_{\Hsch}(s)\le P^f_{\Hsch}(s)$. \textup{(U2c)} Under $\Lone$, quantifier-free $\varphi$ and any prior on $c$, $O_n/O_0\le2^{-n}$.
\end{proposition}

\begin{remark}[$\wedge$-detours]\label{rem:univ:detour}\status{refuted (conjecture); counterexample: proof sketch; computed}\src{universal \S2.3}
The first version conjectured that $\wedge$-detours cannot break the factor $c$ of \Cref{prop:univ:sim}. They can: in $\Cand$ with $(p_{\mathrm{cite}},c,a,e)=(0.35,0.15,0.25,0.25)$ and $\varphi=0+x=x$, $\mu_{\Hall}(\varphi(t))/\mu_{\Hsch}(\varphi(t))=0.15553>c$ (Catalan spine formula, \Cref{app:univ:detour}; samplers $0.15528\pm0.00067$ and $0.1562\pm0.0005$). On a $9139$-point grid the ratio exceeds $c$ but stays $\le0.925$ (normalised $\le0.48$). The factor $c$ is specific to $\Cmin$ and $\CUiii$; whether detours can reverse the direction is open.
\end{remark}

\begin{proposition}[Learned weights under $\Lsel$]\label{prop:univ:B2}\status{proved; computed}\src{universal Prop B2}
In $\Cmin$ let $B=\{\sigma_\psi\}$, $\psi$ quantifier-free, $\inst(\sigma_\psi)\cap\inst(\sigma_\varphi)=\emptyset$, and compare $B\oplus_w\forall x\varphi$ with $B\oplus_w\sigma_\varphi$ (uniform prior on $w$, equal prior masses) under $\Lsel$, $S$ the closed quantifier-free sentences. On data i.i.d.\ from $(1-w^*)\Qg_{\sigma_\psi}+w^*\Qg_{\sigma_\varphi}$ the Bayes factor tends a.s.\ to $h(w^*)$, $h(v):=c/(c+v(1-c))^2\in[c,1/c]$, which exceeds $1$ iff $w^*<\sqrt c/(1+\sqrt c)$.
\end{proposition}

\noindent The cause is the induced prior on the effective weight $wc/(1-w+wc)$ (computed: within $4\cdot10^{-4}$ of $h(w^*)$ at $n=10^5$). The effect is bounded and is not evidence that the axioms prove $\forall x\varphi$.

\begin{remark}[Per-citation selection]\label{rem:univ:B2cit}\status{proved}\src{editor, from experiments \S1.7, Prop X6(b) and universal Prop B2}
Under the per-citation $\Lselc$ both theories of \Cref{prop:univ:B2} have law $(1-w)\Qg_{\sigma_\psi}+w\Qg_{\sigma_\varphi}$ at every $w$, since the selected law of the component $\forall x\varphi$ is $\Qg_{\sigma_\varphi}$ (experiments Prop X6(b)): the Bayes factor is $1$ at every $n$ (cf.\ \Cref{rem:model:sel}).
\end{remark}

\begin{proposition}[Spare-slot prover]\label{prop:univ:both}\status{proved; computed}\src{universal \S4}
In $\{\Hsch,\Hall,\Hboth\}$ with a uniform prior on $w_\forall$, under $\Lone$ with quantifier-free $\varphi$, the likelihood ratio $\Hboth:\Hsch$ is $I_n=\int_0^1(\frac{1-w(1-c)}{1+wc})^ndw\in[\frac1{n+1},\frac1{n(1-c)}]$, so $P(T\vdash\forall x\varphi\mid D_n)\to0$ like $\Theta(1/n)$: a spare slot (\Cref{prop:ident:spare}).
\end{proposition}

The experiments reproduce \Cref{thm:univ:odds} in their chain $\Cch{1}$, whose factor $c_{\mathrm{ch}}=\tfrac12$ matches the unnormalised row: one bit per datum to $1.1\cdot10^{-11}$, and under $\Lselc$ an exact tie with $P(T\vdash\forall x\varphi\mid D_n)\to0.484$ ($0.269$ for $\neg Sx=0$), the prior share in their code (E1, \Cref{tab:exp:e1}).

% ---------------------------------------------------------------------
\subsection{Memorisers, over-general and over-specific templates}\label{sec:univ:memo}

\begin{theorem}[Size principle on instances]\label{thm:univ:size}\status{proved; computed}\src{universal Thm U4}
Let the data be i.i.d.\ from $P^*$, $P^*(I)=1$ for $I=\Ic(\varphi)$, and $P'$ a sub-probability with $P'(I)=p<1$. Then $\KL(P^*\|P')\ge\ln\frac1p$, $\frac1n\ln\frac{P'(D_n)}{P^*(D_n)}\to-\KL(P^*\|P')$ a.s., and $\Ex\frac{P'(D_n)}{P^*(D_n)}\le p^n$. A first-order template strictly more general than $\sigma_\varphi$, with metavariables i.i.d.\ from $\Qg$ and a formula law $F$, has $P_\tau(I)\le\max(\max_fp_f,\max_gF(\mathrm{root}=g))$.
\end{theorem}

\begin{proposition}[Over-specific templates]\label{prop:univ:overspec}\status{proved; computed}\src{universal Prop U5}
For $\tau=\sigma_\varphi[z:=f(z_1,\dots,z_r)]$ and root-generated $\Qg$, $P_\tau(\varphi(t))/P_{\Hsch}(\varphi(t))$ is $1/p_f$ if $\mathrm{root}(t)=f$ and $0$ otherwise. So the odds $\tau:\sigma_\varphi$ grow like $p_f^{-n}$ while all data have root $f$ (probability $p_f^n$), then vanish; they form a martingale. Some single-root template survives with probability $\sum_fp_f^n$, the failure probability of AS Thm~4.4(a).
\end{proposition}

\begin{proposition}[Memorisers]\label{prop:univ:memo}\status{(i), (ii), (ii$'$) proved; (iii) conjecture; computed}\src{universal Prop U6}
Let $Z_0:=\prod_s(1-r_s)>0$, $S_n$ the distinct data, $m_n=\lvert S_n\rvert$, $\mu$ the prior of the memoriser class. \textup{(i)} The class marginal lies in $[Z_0,1]\cdot\prod_{s\in S_n}r_s\,m_n^{-n}$ for uniform weights ($\DirMult_\alpha(D_n)$ in place of $m_n^{-n}$ for Dirichlet weights). \textup{(ii)} Uniform weights, $\Qg$ of infinite support and finite entropy: the log-odds against $\Hsch$ divided by $n$ tend to $-\infty$ a.s. \textup{(ii$'$)} Laplace weights: deterministically the log-odds are $\ge\ln\frac{\mu Z_0}{\pi(\Hsch)}+\sum_{s\in S_n}\ln r_s-(m_n-1)\ln(n+m_n-1)$, which under $\Qnum$ (numeral universe) is $\ge-C\ln^2n$ a.s.\ for all large $n$: the class is not exponentially small. \textup{(iii)} Conjecture: $-\Theta(\ln^2n)$ under $\Qnum$.
\end{proposition}

\noindent So the brief's H4(a), ``mass leaves memorisation and over-general templates exponentially'', is proved for over-general templates and refuted for memorisation in general; under $\QGW$ the computed rate is exponential, as distinct data grow linearly (\Cref{app:univ:memo}). E1 tests only the single memoriser $\Mem(D_n)$; every fixed memoriser dies under well-specification (\Cref{rem:ident:memo}).

% ---------------------------------------------------------------------
\subsection{Predictive confirmation}\label{sec:univ:confirm}

\begin{theorem}[Predictive confirmation]\label{thm:univ:confirm}\status{proved; computed}\src{universal Thm U8}
Let $M(\cdot\mid D_n)=\sum_T\pi_n(T)P_T$, $G:=\{T:P_T(I)=1\}$, $T^*\in G$. \textup{(1)} $\sum_{n\ge0}\Ex[1-M(I\mid D_n)]\le\ln\frac1{\pi(T^*)}$. \textup{(2)} The probability that all future data lie in $I$ is $\pi_n(G)$, and $\pi_n(G)\to1$ a.s. \textup{(3)} $\Ex[1-\pi_n(G)]\le\min\big(1,\sum_{T\notin G}\frac{\pi(T)}{\pi(T^*)}P_T(\operatorname{supp}P_{T^*})^n\big)$. \textup{(4)} If $T\in G$ have law $\hat P$ and $T\notin G$ have law $(1-\varepsilon_T)\hat P+\varepsilon_TN$, $N(I)=0$, then on instance sequences $1-\pi_n(G)=\frac{\sum_{T\notin G}\pi(T)(1-\varepsilon_T)^n}{\pi(G)+\sum_{T\notin G}\pi(T)(1-\varepsilon_T)^n}$.
\end{theorem}

\noindent The rate depends on near-universal competitors: $\approx(1-\pi_0)/(\pi_0n)$ for Laplace's rule with a point mass $\pi_0$, but $1-\pi_n(G)=0.37$ at $n=10^{50}$ for a dyadic family ($\varepsilon_j=2^{-j}$, $\pi_0=0.01$).

\begin{remark}[Relation to Hutter]\label{rem:univ:hutter}\status{known; computed}\src{universal \S5}
In the Bayes--Laplace model the universal hypothesis, even its observable form, has posterior $0$, while the universal prior confirms it \citep{hutter2007universal} (abstract checked; full text not). \Cref{thm:univ:confirm}(2) is the analogue for countable classes, by the same mechanism: prior mass on instance-only hypotheses. But $G$ contains $\Hsch$, which does not prove $\forall x\varphi$. Nor is confirmation stepwise monotone: with prior $\tfrac12$ on each of $\sigma_\varphi$, $\varphi(Sz)$, the datum $\varphi(S0)$ lowers $\sigma_\varphi$ to $\tfrac13$, a Nicod-type effect \citep{leike2015solomonoff}.
\end{remark}

% ---------------------------------------------------------------------
\subsection{The Bayesian \texorpdfstring{$\omega$}{omega}-gap}\label{sec:univ:omega}

\begin{theorem}[Bayesian $\omega$-gap]\label{thm:univ:omega}\status{(a), (b), (c1)--(c3), (d1), (e) proved; (d2) proved for the pair, proof sketch in general; computed}\src{universal Thm U10}
Let $\Hclass$ be countable with a proper prior $\pi>0$, $\ProvAll:=\{T\in\Hclass:T\vdash\forall x\varphi\}$ (first-order; inconsistent theories included), and the data i.i.d.\ from $P_{T^*}$, $T^*\in\Hclass$, unless stated.
\begin{enumerate}[label=(\alph*),leftmargin=*,itemsep=1pt]
\item $P(T\vdash\forall x\varphi\mid D_n)\to\pi(\ProvAll\cap\Cstar)/\pi(\Cstar)$ a.s.; likewise for $\vdash_{\mathsf C}$ and for $\Bel,\Dis,\Indep,\Inc$ (\Cref{def:sound:tri}).
\item \emph{Prior share.} If $\Cstar$ has a prover and a non-prover, the limit is in $(0,1)$: $\Hall,\Hsch$ under $\Lcl$ or $\Lsel$; $B\oplus\forall x\varphi$, $B\oplus\sigma_\varphi$ under $\Lcl$ with a common weight or weight prior. In a class of the two theories of \Cref{prop:univ:B2} under $\Lsel$ the laws differ, but that proposition gives a limit in $(0,1)$ directly.
\item \emph{Faithful regime.} Call a likelihood faithful for $\mathsf C$ if $\operatorname{supp}P_T=\Th_{\mathsf C}(T)\cap X$ for all $T$ ($X$ the data space). (c1) If $\forall x\varphi\in X$, $P(T\vdash_{\mathsf C}\forall x\varphi\mid D_n)\to1[T^*\vdash_{\mathsf C}\forall x\varphi]$. (c2) If also $\mathsf C$ is sound and all axiom instances lie in $X$, members of $\Cstar$ have $T^*$'s first-order theorems and $P(T\vdash\forall x\varphi\mid D_n)\to1[T^*\vdash\forall x\varphi]$. (c3) In $\Cmin$, and $\Copen$ with a guard, $\Lone$ with full-support $\Qg$ is faithful and $\operatorname{supp}P_{\Hsch}=\Ic(\varphi)$ for quantifier-free $\varphi$, so for $T^*=\Hsch$ the limit in (c1) is $0$. If $\mathsf C$ has a rule that maps an instance to a non-instance with positive probability ($\wedge$I, $\vee$I, symmetry of $=$, $\to$I, MP with cited logical axioms), every $T$ with $P_T(\Ic)>0$ emits non-instances: instance-only data are then misspecified for every hypothesis, and the limit is set by the best-fitting generators of the class (\Cref{thm:ident:kl}).
\item \emph{Filtered data} from $\Hall$, seen only in $S=\Ic(\varphi)$. (d1) For quantifier-free $\varphi$ in $\Cmin$ the filtered law is $\Hsch$'s: limit $0$ under $\Lone$, the prior share under $\Lsel$. (d2) For quantified $\varphi$, $\Hsch$ is the KL-minimiser in $\{\Hall,\Hsch\}$ and the limit is $0$; for larger classes the minimiser was not determined.
\item \emph{Unknown filter.} With a countable filter family $\mathcal F$, prior $\nu$ independent of $\pi$, and pair laws $P_T(s)1[s\in S]/P_T(S)$, (a) applies to pairs. For $\{\Hall,\Hsch\}$, $\Cmin$, quantifier-free $\varphi$ and data law $\Qg_{\sigma_\varphi}$ the limit is $\frac{\pi(\Hall)\nu(\mathcal F_1)}{\pi(\Hall)\nu(\mathcal F_1)+\pi(\Hsch)\nu(\mathcal F_0)}$, $\mathcal F_0=\{S\supseteq\Ic\}$, $\mathcal F_1=\{S\in\mathcal F_0:\forall x\varphi\notin S\}$.
\end{enumerate}
\end{theorem}

\noindent\emph{Proof idea.} (a) is \Cref{thm:ident:doob} summed over $\ProvAll$; the rest identifies $\Cstar$ in each setting (\Cref{app:univ:omega}). Computed: on closed data in $\Copen$ the prover mass falls from $0.750$ to $0.004$ at $n=50$ under $\Lone$ and stays $0.750$ under $\Lsel$ (\Cref{tab:univ:c5}); (e) reaches its predicted $0.2500$ by $n=10$.

In the prior-share regime the answer is the prior share. That it lies in $(0,1)$ is robust; its value is not: for $\varphi=0+x=x$ the share of $\Hall$ in $\{\Hall,\Hsch\}$ is $0.333$ in the universal code and in $\pi_1$ of \Cref{def:model:prior}, $0.200$ in $\pi_2$, $0.484$ in the experiments' code and $0.030$--$0.042$ in pa's codes (\Cref{tab:univ:share}). The data never move it.

\begin{proposition}[H\"anni's variants on $\forall x\varphi$]\label{prop:univ:hanni}\status{proved; computed}\src{universal \S6.2}
Let $\varphi=0+x=x$, $\Qg$ of full support. Under $\Sprove$ the posterior converges a.s.\ in total variation to the prior restricted to $\{T\text{ consistent}:T\vdash\Ic(\varphi)\}$, which contains $\Q$ (independent), $\Q+\forall x\varphi$ (true), $\Q+\exists x\neg\varphi$ (false), $\Hsch$, $\Hall$ and consistent over-general templates such as $\{0+z_1=z_2\}$: all three trichotomy values keep positive mass, and over-generalisation is never penalised. The same holds for $\Snc$ and $\Sg$ ($g\equiv1$ on provable data). Under generative likelihoods inconsistent theories die within a few data (\Cref{tab:univ:inc}).
\end{proposition}

\noindent Inconsistent theories inflate $\Bel$ and $\Dis$ alike (at $n=0$ in the class of \Cref{tab:univ:c2}, $\Bel(\forall x\varphi)\approx0.69$ is mostly $\Inc$), so the verifier should use $\Belc$ (\Cref{rem:sound:belc}).

\begin{proposition}[An $\omega$-gap for Q4]\label{prop:univ:q4}\status{proved; computed}\src{universal Prop U13}
$(\Q-\text{Q4})$ plus all closed instances of $x+0=x$ does not prove $\forall x(x+0=x)$ (a model on $\N\cup\{a,b\}$). So closed instances cannot tell the axiom Q4 from its closed schema $z+0=z$, which differ in their theorems over $\Q-\text{Q4}$.
\end{proposition}

\begin{proposition}[The Gaifman condition]\label{prop:univ:gaifman}\status{proved; known}\src{universal Prop U15}
Let $\mu$ be a probability on $L$-structures, $P(\psi):=\mu\{M:M\models\psi\}$, and $t_1,t_2,\dots$ the closed terms. Then
\[P\Big(\forall x\varphi\ \Big|\ \textstyle\bigwedge_{i\le n}\varphi(t_i)\Big)\to P(\forall x\varphi)\big/\mu\{M:M\models\varphi(t)\text{ for all closed }t\}.\]
The Gaifman condition \citep{gaifman1964concerning} $P(\forall x\varphi)=\inf_nP(\bigwedge_{i\le n}\varphi(t_i))$ holds for all $\varphi$ iff $M_0\preccurlyeq M$ for $\mu$-a.e.\ $M$, with $M_0$ the closed-term substructure (AS Prop~4.7); the limit is then $1$ if $P(\forall x\varphi)>0$ and $0$ if $P(\forall x\varphi)=0$.
\end{proposition}

\noindent So believing $\forall x\varphi$ from true closed instances is safe for $\N$ and unsafe for credences that weigh $\R$, $V$, models of $\PA+\neg\Con(\PA)$ (AS Ex~4.8) or some models of $\Q$ (AS Ex~4.9). The condition is cited from memory by the track; secondary sources state it with constants rather than closed terms.

% ---------------------------------------------------------------------
\subsection{Open instances and the closedness guard}\label{sec:univ:open}

\begin{proposition}[Parameter data and the guard]\label{prop:univ:open}\status{proved; computed}\src{universal Prop U11(a)--(d)}
Assume the closure reading, Gen over parameters, theories without parameters, and a likelihood supported on theorems. \textup{(a)} Once $D$ contains some $\varphi(p)$, every $T$ with $\pi(T\mid D)>0$ proves $\forall x\varphi$. \textup{(b)} If $P_{T^*}(\varphi(p))=\beta>0$, then $P(P(T\vdash\forall x\varphi\mid D_n)<1)\le(1-\beta)^n$. \textup{(c)} In $\Copen(c,g)$ with $\Qopen$, per closed datum under $\Lone$: $\Hopen:\Hsch=\frac{1-\rho}{1+g\rho}$, $\Hall:\Hsch=\frac{c(1-\rho)}{1+c}$, $\Hall:\Hopen=\frac{c(1+g\rho)}{1+c}$. \textup{(d)} If the guarded $\Hsch$ is in the class, closed data give $P(T\vdash\forall x\varphi\mid D_n)\to0$: the guard is learned at rate $\Wo:=\ln\frac{1+g\rho}{1-\rho}$ ($0.1252$ at the defaults), and $\Hboth$ decays like $1/n$.
\end{proposition}

\noindent Here $\rho$ is fixed. In pa's two-part code, where the parameter probability is learned, the analogous head-symbol split of $x+0=x$ closes in on the schema at only about 5 bits per doubling of $n$ (\Cref{rem:pa:merge}); each rate holds in its model.

\begin{remark}[Without a guard]\label{rem:univ:openrefuted}\status{refuted; computed}\src{universal \S7.2}
The first version claimed: ``without a guard ($\DT$, $\lambda$-convention, parameters admissible) $\Hopen$ wins on closed data and every surviving hypothesis proves $\forall x\varphi$: the posterior takes the $\omega$-step''. Counterexample: the unguarded $\DT$ pair $R_1=\{\varphi(0),\varphi(Sz)\}$ beats $\Hopen$ by $(1-q)\Wo$ per datum (\Cref{prop:univ:rkrate}) and does not prove $\forall x\varphi$ in pure logic; in $\{\Hall,\Hopen,R_1\}$, $P(T\vdash\forall x\varphi\mid D_n)$ is $0.697$ at $n=100$ and $0.000$ from $n=1000$ (referee: $0.880$, then $0.000$ at $300$). The claim survives for the class $\{\Hall,\Hopen\}$ and under $\Lsel$. This is the Bayesian counterpart of AS Prop~4.5(b), (c).
\end{remark}

\begin{proposition}[Rates on closed data]\label{prop:univ:rkrate}\status{proved; computed}\src{universal Prop N2}
On closed numeral data from $\Qnum$, under $\Lone$: $\KL(P^*\|P_{\Hall})=\ln\frac{1+c}{c(1-\rho)}$, $\KL(P^*\|P_{\Hopen})=\Wo$, $\inf_w\KL(P^*\|P_{R_k})=q^k\Wo$. Each split level halves the loss at $q=\tfrac12$ ($0.0626$, $0.0313$, $0.0156$ for $R_1,R_2,R_3$).
\end{proposition}

\noindent The proof uses the exact law of $R_k$ (\Cref{prop:univ:rk}) and a waste lemma (\Cref{lem:univ:waste}): a law that reproduces $P^*$ on a tail of mass $\tau$ and wastes a fixed fraction $r_{\mathrm w}$ of the tail's mass loses at best $\tau\ln(1+r_{\mathrm w})$; for $R_k$, $\tau=q^k$ and $\ln(1+r_{\mathrm w})=\Wo$.

\begin{proposition}[Who wins without a guard]\label{prop:univ:noguard}\status{provability, (a), (b) proved; (c) proof sketch; computed; (d) computed, limit conjecture}\src{universal Prop N3}
Closed data i.i.d.\ from $\Qnum$; the class contains $\Hall$, $\Hopen$, $R_k$ for $k$ in a nonempty $\mathcal K$ (learned or optimal weights), possibly memorisers. In pure logic $\Hall$, $\Hopen$ prove $\forall x\varphi$, and $R_k$ and memorisers do not; over $\Q$ all but the memorisers do. \textup{(a)} $\Lone$, pure logic, provers only $\Hall,\Hopen$: $P(T\vdash\forall x\varphi\mid D_n)\to0$ a.s., exponentially, at rate $\ge(1-q^k)\Wo$ for each $k\in\mathcal K$. \textup{(b)} $\Lone$, over $\Q$, $\mathcal K$ finite: $\to0$ with Laplace-weighted memorisers, $=1$ at every $n$ without them. \textup{(c)} $\Lsel$, finite $\mathcal K$: $\to1$ in pure logic. \textup{(d)} $\Lone$, $\mathcal K=\{1,\dots,40\}$, memorisers: $1.000$ over $\Q$ from $n=10$ to $10^{12}$, with ever deeper splits.
\end{proposition}

\noindent\emph{Proof idea.} $R_k$ fails $\forall x\varphi$ in $\N\cup\{e,e'\}$ with $Se=Se'=e'$, $0+e=0+e'=e'$; over $\Q$, $k$ uses of Q3 reduce $\forall x\varphi$ to closed instances and $\forall y\varphi(S^ky)$. (a), (b) compare the linear rates of \Cref{prop:univ:rkrate} with each other and with the $O(\ln^2n)$ cost of memorisers; (c) uses the $\frac k2\ln n$ Occam term \citep{schwarz1978estimating,clarke1990information}, the almost-sure uniform step not written out.

In c8 (\Cref{tab:univ:noguard}), with $R_{1..4}$ and memorisers the over-$\Q$ value is $1.000$ up to $n=10^4$ and $0.000$ from $10^5$; with $R_{1..40}$ it stays $1.000$ to $n=10^{12}$, the MAP moving from $R_3$ ($n=10^3$) to $R_{30}$ ($10^{12}$). So without a guard, under a likelihood that ignores the selection, the posterior does not take the $\omega$-step in pure logic; it approximates the missing guard by ever deeper splits, whose waste $q^k\Wo$ tends to $0$. Over $\Q$ the answer is a race between the split family (provers) and the memoriser family (non-provers): a property of the class. Experiments E3(a) show a related drift when the term law, not a filter, is misspecified: the best closed-guard split $R_m$ deepens from $m=9$--$10$ at $n=512$ to $11$--$12$ at $1024$ (\Cref{tab:exp:e3a}); these runs are too short to test (b).

% ---------------------------------------------------------------------
\subsection{Quantified data, lemmas and the equational fragment}\label{sec:univ:quant}

\begin{proposition}[Quantified data]\label{prop:univ:quant}\status{(a) proved; (b) $f_q=0$ proved, $f_q>0$ proof sketch; computed; (c) proof sketch}\src{universal Prop U12}
\textup{(a)} If $s\in D$ and $B\cup\Ic(\varphi)\nvdash s$, then under any likelihood supported on theorems $B\oplus\sigma_\varphi$ has posterior $0$; over $\Q$, one datum $\forall y(0+Sy=Sy)$ leaves only provers of $\forall x(0+x=x)$. \textup{(b)} If all theories prove $\forall x\varphi$ at different costs ($\{\forall^L\forall x\varphi,\sigma_\varphi\}$ against $\{\forall x\varphi,\sigma_\varphi\}$ in $\Cmin$, $L$ vacuous quantifiers, uniform weight priors, equal prior masses, data $\forall x\varphi$ with probability $f_q$ and instances otherwise), the log Bayes factor tends to $\ln(1+c+\dots+c^L)$ when $f_q=0$ ($0.349$ at $L=3$) and grows linearly when $f_q>0$, at a rate a little above $f_qL\ln\frac1c$. \textup{(c)} In $\Lmax$ form a derived sentence becomes an axiom when the derivation steps it saves, at $\ge\ln\frac1c$ each, exceed its prior cost and an $O(\ln n)$ dilution: frequently used lemmas become axioms.
\end{proposition}

\noindent So the posterior identifies the data producer's working set of primitives, not a canonical axiomatisation (compare \Cref{prop:pa:usage}). The brief's ``statements given without proof should be easily derivable'' was tested on the equational fragment $E=\{\text{Q4},\text{Q5}\}$ (\Cref{prop:univ:eqfrag}): every derivation of $0+S^k0=S^k0$ has at least $k+1$ axiom leaves, so the two-part ($\Lmax$) comparison favours the schema $0+z=z$ linearly in $k$ at $q=\tfrac12$, while the full-sum comparison is shown to do so only under the sufficient condition $q>e^{-(1-\mu)^2/2}$ ($\mu$ the mean number of premises), which needs $q>0.61$ and says nothing at $q=\tfrac12$. Full $\Q$, and the full sum at $q=\tfrac12$, are open. A Kt-style penalty on generating $0+S^k0=S^k0$ from $E$ would charge only $O(\ln k)$ (\Cref{rem:time:links}).

% ---------------------------------------------------------------------
\subsection{Without templates}\label{sec:univ:sentences}

H\"anni proposes templates against his collapse. Without them, hypotheses are finite sets of sentences and $\sigma_\varphi$ is unavailable. Setting: $\Cmin$, $c=0.3$, numeral data, $q=\tfrac12$, $\varphi=0+x=x$, learned weights.

\begin{proposition}[Sentence-only class]\label{prop:univ:sentences}\status{(a) proof sketch; computed; (b), (c) proved; (d) computed, limit conjecture}\src{universal Prop U16}
\textup{(a)} Under $\Lcl$, $\Hall$ has the true law and $\Split_k$, Both0 match it only at a point or at the boundary, so $P(T\vdash\forall x\varphi\mid D_n)\to1$ (Laplace expansion uniform in $k$ not written out). \textup{(b)} Under $\Lone$: $\KL(P^*\|P_{\Hall})=\ln\frac{1+c}c$ ($1.466$), $\inf_w\KL(P^*\|P_{\Split_k})=q^k\ln\frac{1+c}c$, $\inf_w\KL(P^*\|P_{\mathrm{Both0}})=q\ln\frac{1+qc}{qc}$ ($1.018$). \textup{(c)} $\Split_k$ proves $\forall x\varphi$ over $\Q$, not in pure logic; memorisers in neither. \textup{(d)} In pure logic $P(T\vdash\forall x\varphi\mid D_n)\to0$ in every class computed. Over $\Q$, with depth $\le4$ memorisers win (limit $0$); with depth $\le40$ splits win to $n=10^{12}$ under four prior variants (\Cref{tab:univ:sentences}). Conjecture: with all depths the limit is $1$.
\end{proposition}

\begin{lemma}[Survivors prove $\forall x\varphi$ over $\Q$]\label{lem:univ:survivors}\status{proved}\src{universal Lemma S1}
In $\Cmin$ with numeral elimination terms and $\varphi=0+x=x$, every finite set $T$ of sentences with $P_T(\varphi(S^j0))>0$ for infinitely many $j$ satisfies $\Q\cup T\vdash\forall x\varphi$.
\end{lemma}

\noindent A fixed memoriser is eventually refuted but the memoriser family is not, so the over-$\Q$ limit is a race between families. The first version's ``$=1$ over $\Q$'' (a three-hypothesis class) and its referee's ``$\to0$ over $\Q$'' (bounded depth) were class artefacts (\Cref{app:univ:sentences}).

\begin{remark}[Kreisel's conjecture]\label{rem:univ:kreisel}\status{known (secondary sources only)}\src{universal \S9}
``All instances have short derivations, so $\forall x\varphi$ is derivable'' resembles Kreisel's conjecture, reported open for standard formalisations of $\PA$, settled for some \citep{parikh1973some,baaz1993kreisel} and false for some theories close to $\PA$ \citep{hrubes2007theories} (none of these was opened). $\Split_1$ shows that the naive version fails without axioms such as Q3: every numeral instance has a one-step derivation, $\forall x\varphi$ none.
\end{remark}

% ---------------------------------------------------------------------
\subsection{The intuition, precisely}\label{sec:univ:intuition}

The intuition is right against hypotheses that fit worse (over-general templates, memorisers), for prediction (``all future data are instances''), and for truth under a Gaifman-type credence, safely in $\N$. It needs refinement in three places. $\sigma_\varphi$ is a second simple law that implies the data, and it is never overtaken (\Cref{thm:univ:B}); templates, proposed against the collapse, are what bring it into the class. ``The axioms prove $\forall x\varphi$'' does not gain from closed instances (\Cref{thm:univ:omega}). A derivation likelihood rewards the most specific generator that fits, and frequently used lemmas, not the logically strongest law. What licenses $\forall x\varphi$ is an instance at a parameter, a quantified consequence, or a selection-aware likelihood without a guard (\Cref{prop:univ:open}(a), \Cref{prop:univ:quant}(a), \Cref{prop:univ:noguard}(c)), which takes the $\omega$-step by design, soundly in $\N$ and unsoundly in $\R$. Data at parameters are common in proof-assistant practice, so the practical answer is often positive, for a reason other than the intuition's.

A good version for this case: a proper template prior, inconsistent theories excluded or reported as $\Inc$; a derivation likelihood that models the selection of reported data, by a known filter ($\Lsel$) or a prior over filters (\Cref{thm:univ:omega}(e)); a noise component (\Cref{prop:univ:robust}); and the verifier ``accept $s$ iff $\Belc(s)\ge1-\delta$''. On closed instances it accepts $\forall x\varphi$ only if the prior share of provers over theories and filters is $\ge1-\delta$; one parameter instance, or one quantified datum entailing it over the background, suffices.

\begin{corollary}[Soundness for $\forall x\varphi$]\label{cor:univ:sound}\status{proved}\src{universal \S11}
Let the data be i.i.d.\ from $P_{T^*}$ with \emph{fixed} weights, $T^*\in\Hclass$, $T^*\nvdash\forall x\varphi$, and $C^*_\forall:=\{T\in\Cstar:T\vdash\forall x\varphi\Leftrightarrow T^*\vdash\forall x\varphi\}$. The probability that the verifier (on $\Bel$ or $\Belc$) ever accepts $\forall x\varphi$ is at most $\delta/\pi(C^*_\forall)\le\delta/\pi(T^*)$, uniformly in time: \Cref{thm:sound:fixed} read at the sentence $\forall x\varphi$.
\end{corollary}

\noindent\emph{Proof.} Members of $C^*_\forall$ do not prove $\forall x\varphi$ and share $T^*$'s likelihood, so
\[\pi_n(C^*_\forall)=\pi(C^*_\forall)\,P_{T^*}(D_n)/M(D_n).\]
Acceptance needs $\pi_n(C^*_\forall)\le\delta$, that is, $M(D_n)/P_{T^*}(D_n)\ge\pi(C^*_\forall)/\delta$. This ratio is a nonnegative supermartingale started at $1$; by Ville's inequality it ever reaches that level with probability at most $\delta/\pi(C^*_\forall)$. \qed

Without well-specification there is no guarantee. In the misspecified settings of \Cref{thm:univ:omega}(d) and of \Cref{prop:univ:noguard,prop:univ:sentences}, the limit theories ($\sigma_\varphi$, $R_k$, $\Split_k$, memorisers) are true in $\N$, so in pure logic the verifier errs towards incompleteness; misspecification can also make it accept a falsehood with probability $1$ (\Cref{ex:ident:escape,prop:sound:misspec}). The learned-weight computations here are outside the corollary; for them see \Cref{thm:sound:avg,thm:sound:shrink}.

\begin{remark}[Relation to the MDL finding]\label{rem:univ:mdl}\status{proved (from the cited results)}\src{universal \S10}
AS \S6.10 found that ``MDL tracks the statistics of usage, not the logical boundaries of schemas''. For $\forall x\varphi$: well specified, a root split with weights $p_f$ has exactly the schema's law under $\Lzero$ and $\Lone$ (\Cref{prop:ident:splits}; it keeps its prior share $0.004$ in \Cref{tab:univ:c2}), and with learned weights it pays a polynomial Occam factor (\Cref{prop:ident:splitlzero}; AS Prop~E.17); an unmodelled selection creates linear split wins, at rates $(1-q^k)\Wo$ for $R_k$ and $(1-q^k)\ln\frac{1+c}c$ for $\Split_k$ (the analogue of AS Prop~E.16, split by numeral depth rather than main connective); modelling the selection removes them (\Cref{prop:univ:noguard}(c)). So a derivation likelihood does not change the MDL finding; it adds a way to be misspecified, since the theorems humans state are a selection from those their axioms generate (\Cref{rem:ident:mdl}).
\end{remark}
